% Generated draft: requires mathematical and source review.
\documentclass[11pt]{article}
\usepackage[margin=25mm]{geometry}
\usepackage{fontspec}
\setmainfont{Latin Modern Roman}
\usepackage{amsmath,amssymb,mathtools,unicode-math}
\setmathfont{Latin Modern Math}
\usepackage{enumitem,needspace,xurl,hyperref}
\setlength{\parindent}{0pt}
\setlength{\parskip}{5pt}
\setlength{\emergencystretch}{3em}
\newcommand{\N}{\mathbb{N}}
\newcommand{\Z}{\mathbb{Z}}
\newcommand{\Q}{\mathbb{Q}}
\newcommand{\R}{\mathbb{R}}
\newcommand{\C}{\mathbb{C}}
\newcommand{\F}{\mathbb{F}}
\newcommand{\A}{\mathbb{A}}
\newcommand{\PP}{\mathbb{P}}
\newcommand{\E}{\mathbb{E}}
\newcommand{\eps}{\varepsilon}
\newcommand{\dd}{\,\mathrm{d}}
\begin{document}
\section*{UnsolvedMath ID: 701192}
\textbf{Source problem code:} catalogue entry
\subsection{Definitions and mathematical statement}
A group word is an element \(w\) of a free group on variables \(x_1,x_2,\ldots\); each word involves only finitely many variables. Choose \(m\geq 1\) such that \(w\) involves only \(x_1,\ldots,x_m\). For a group \(G\), a \(w\)-value is an element \(w(g_1,\ldots,g_m)\) obtained by substituting \(g_i\in G\). The verbal subgroup associated with \(w\) is
\[
 w(G)=\left\langle w(g_1,\ldots,g_m):g_1,\ldots,g_m\in G\right\rangle.
\]
The identity \(w^n=1\) means that
\[
 w(g_1,\ldots,g_m)^n=1_G
 \qquad\text{for every }(g_1,\ldots,g_m)\in G^m.
\]
A group \(G\) is residually finite if, for every \(g\neq 1_G\), there exist a finite group \(Q\) and a homomorphism \(\varphi:G\to Q\) such that \(\varphi(g)\neq 1_Q\). A group is locally finite if each of its finitely generated subgroups is finite.

\textbf{Question.} For every positive integer \(n\), every group word \(w\), and every residually finite group \(G\) satisfying \(w^n=1\), must \(w(G)\) be locally finite? Equivalently, under these hypotheses, must
\[
 \bigl|\langle h_1,\ldots,h_r\rangle\bigr|<\infty
 \qquad\text{for every }r\geq 1\text{ and all }h_1,\ldots,h_r\in w(G)?
\]
No finite generation assumption is imposed on \(G\).

The supplied source identifies the specialization \(w=x_1\), for which \(w(G)=G\), with the Restricted Burnside Problem. It also reports positive answers in a number of other particular cases, citing the article by P. V. Shumyatsky listed below; those cases are not specified in the supplied record. The general question is labeled open in that record.
\subsection{Short English statement}
Must the verbal subgroup generated by the values of any group word be locally finite whenever the ambient group is residually finite and all values of that word have orders dividing one fixed positive integer? This is Problem 15.104, proposed by P. V. Shumyatsky in 2002; its open status is supplied registry metadata.
\subsection{Sources}
\begin{itemize}
\item Supplied local UnsolvedMath registry record, permanent ID \(701192\), exact source code \texttt{catalogue entry}, titled \emph{Kourovka Notebook Problem 15.104}. The record is marked open and was updated on 2026-10-16.
\item \emph{The Kourovka Notebook (Unsolved Problems in Group Theory)}, original Problem 15.104, 15th issue (2002); the supplied background also identifies the 21st issue edition (2026), \texttt{arXiv:1401.0300v48}. Proposed by P. V. Shumyatsky. The supplied collection description names E. I. Khukhro and V. D. Mazurov as editors.
\item P. V. Shumyatsky, \emph{Quart. J. Math.}, volume \(51\), no.\ \(4\) (2000), pp.\ 523--528. The supplied statement cites this article for positive answers in a number of particular cases; it supplies neither the article title nor the precise cases.
\end{itemize}
\end{document}
